% GENERATED FILE. Edit canonical lessons, then run npm run generate:books.
% canonical-source-hash: fnv1a64:9bec12e7880d5c6e
% canonical-lessons: JA-C103-otama, JA-C103-suidou, JA-C103-yakan, JA-C103-obon, JA-C103-kan

\chapter{Ladle, Water supply, Kettle, Tray, Can}
\label{ch:ja-ladle-water-supply-kettle-tray-can}
\hlchaptermodality{103}

\begin{chapteropening}
\textbf{By the end of this chapter:} \emph{I can say a ladle, the water supply, a kettle, a tray and a can.}

Complete the last lesson of chapter 103: i can say a ladle, the water supply, a kettle, a tray and a can.
\end{chapteropening}

\section[\texorpdfstring{\ja{おたま} (otama)}{おたま (otama)}]{\ja{おたま} (otama) — a ladle}

\label{lesson:JA-C103-otama}



\begin{quote}
Before the new one: say the Japanese for a rice pot, then the Japanese for a chopping board.
\end{quote}

\subsection*{You'll want to know: \ja{おたま}}
\textbf{\ja{おたま}} — \emph{otama} — \textquotedblleft{}a ladle\textquotedblright{}.

\begin{grammarlens}[title={What you've built}]
One more word for this chapter.
\end{grammarlens}

\subsection*{Your turn}
\begin{itemize}
\raggedright
  \item \emph{Say it:} \emph{otama}
  \item \emph{Say it:} \emph{otama}, once more
  \item \emph{Say it:} say \emph{otama}
  \item \emph{Recall:} say the Japanese for a rice pot, then the Japanese for a chopping board, then say \emph{otama} again
\end{itemize}

\begin{tcolorbox}[breakable,colback=teal!4,colframe=teal!35!black,title={Before you move on}]
What does \ja{おたま} mean? (\textquotedblleft{}A ladle\textquotedblright{}.) Say it once more.
\end{tcolorbox}
\section[\texorpdfstring{\ja{すいどう} (suidou)}{すいどう (suidou)}]{\ja{すいどう} (suidou) — the water supply}

\label{lesson:JA-C103-suidou}



\begin{quote}
Before the new one: say the Japanese for a chopping board, then the Japanese for a ladle.
\end{quote}

\subsection*{You'll want to know: \ja{すいどう}}
\textbf{\ja{すいどう}} — \emph{suidou} — \textquotedblleft{}the water supply\textquotedblright{}.

\begin{grammarlens}[title={What you've built}]
One more word for this chapter.
\end{grammarlens}

\subsection*{Your turn}
\begin{itemize}
\raggedright
  \item \emph{Say it:} \emph{suidou}
  \item \emph{Say it:} \emph{suidou}, once more
  \item \emph{Say it:} say \emph{suidou}
  \item \emph{Recall:} say the Japanese for a chopping board, then the Japanese for a ladle, then say \emph{suidou} again
\end{itemize}

\begin{tcolorbox}[breakable,colback=teal!4,colframe=teal!35!black,title={Before you move on}]
What does \ja{すいどう} mean? (\textquotedblleft{}The water supply\textquotedblright{}.) Say it once more.
\end{tcolorbox}
\section[\texorpdfstring{\ja{やかん} (yakan)}{やかん (yakan)}]{\ja{やかん} (yakan) — a kettle}

\label{lesson:JA-C103-yakan}



\begin{quote}
Before the new one: say the Japanese for a ladle, then the Japanese for the water supply.
\end{quote}

\subsection*{You'll want to know: \ja{やかん}}
\textbf{\ja{やかん}} — \emph{yakan} — \textquotedblleft{}a kettle\textquotedblright{}.

\begin{grammarlens}[title={What you've built}]
One more word for this chapter.
\end{grammarlens}

\subsection*{Your turn}
\begin{itemize}
\raggedright
  \item \emph{Say it:} \emph{yakan}
  \item \emph{Say it:} \emph{yakan}, once more
  \item \emph{Say it:} say \emph{yakan}
  \item \emph{Recall:} say the Japanese for a ladle, then the Japanese for the water supply, then say \emph{yakan} again
\end{itemize}

\begin{tcolorbox}[breakable,colback=teal!4,colframe=teal!35!black,title={Before you move on}]
What does \ja{やかん} mean? (\textquotedblleft{}A kettle\textquotedblright{}.) Say it once more.
\end{tcolorbox}
\section[\texorpdfstring{\ja{おぼん} (obon)}{おぼん (obon)}]{\ja{おぼん} (obon) — a tray}

\label{lesson:JA-C103-obon}



\begin{quote}
Before the new one: say the Japanese for the water supply, then the Japanese for a kettle.
\end{quote}

\subsection*{You'll want to know: \ja{おぼん}}
\textbf{\ja{おぼん}} — \emph{obon} — \textquotedblleft{}a tray\textquotedblright{}.

\begin{grammarlens}[title={What you've built}]
One more word for this chapter.
\end{grammarlens}

\subsection*{Your turn}
\begin{itemize}
\raggedright
  \item \emph{Say it:} \emph{obon}
  \item \emph{Say it:} \emph{obon}, once more
  \item \emph{Say it:} say \emph{obon}
  \item \emph{Recall:} say the Japanese for the water supply, then the Japanese for a kettle, then say \emph{obon} again
\end{itemize}

\begin{tcolorbox}[breakable,colback=teal!4,colframe=teal!35!black,title={Before you move on}]
What does \ja{おぼん} mean? (\textquotedblleft{}A tray\textquotedblright{}.) Say it once more.
\end{tcolorbox}
\section[\texorpdfstring{\ja{かん} (kan)}{かん (kan)}]{\ja{かん} (kan) — a can, a tin}

\label{lesson:JA-C103-kan}



\begin{quote}
Before the new one: say the Japanese for a kettle, then the Japanese for a tray.
\end{quote}

\subsection*{You'll want to know: \ja{かん}}
\textbf{\ja{かん}} — \emph{kan} — \textquotedblleft{}a can, a tin\textquotedblright{}.

\begin{grammarlens}[title={What you've built}]
One more word for this chapter.
\end{grammarlens}

\subsection*{Your turn}
\begin{itemize}
\raggedright
  \item \emph{Say it:} \emph{kan}
  \item \emph{Say it:} \emph{kan}, once more
  \item \emph{Say it:} say \emph{kan}
  \item \emph{Recall:} say the Japanese for a kettle, then the Japanese for a tray, then say \emph{kan} again
\end{itemize}

\begin{tcolorbox}[breakable,colback=teal!4,colframe=teal!35!black,title={Before you move on}]
What does \ja{かん} mean? (\textquotedblleft{}A can, a tin\textquotedblright{}.) Say it once more.
\end{tcolorbox}
